\honest{The Proper Use of Doubt}{Eliezer Yudkowsky}{2007}

I once remarked that most organized belief systems exist to flee from doubt. A listener replied that the Jesuits must be an exception: their novices, he said, are told to doubt Christianity, God and their own calling, and even to doubt that their doubts may grow stronger. Googling did not confirm or refute this, so I treat it as a possibly hypothetical case.

Such Jesuits, I concede, would not be fleeing from doubt. But the practice strikes me as ``highly suspicious.'' To a truly virtuous rationalist doubt should not be scary, and this sounds like exposing an arachnophobe to spiders. \nb{The post's own last rule grants that ``it may take courage to face your doubts,'' and says only that ``to an ideal mind'' doubt would not be scary. For minds that are not ideal, practice at facing doubt is what that rule would seem to call for.}

Curiosity wants an answer, and so seeks its own end. In the same way, every doubt exists to destroy a particular belief, or to die trying. A rational doubt arises for a specific reason, and that reason points to an investigation, which ends either the belief or the doubt. \nb{Charles Peirce made this argument against Descartes: in 1868, ``Let us not pretend to doubt in philosophy what we do not doubt in our hearts,'' and in 1877, ``There must be a real and living doubt, and without this all discussion is idle.'' The post does not mention him.} ``An unresolved doubt is a null-op; it does not turn the wheel, neither forward nor back.'' \nb{Too strong, by Yudkowsky's own example. In ``The Proper Use of Humility,'' an engineer who is sure a machine will not fail builds fail-safes anyway; that doubt is not resolved, and it changes the design.}

So why would a real believer give novices doubts that must fail? It would be like asking physics students to agonize over whether Newton was right after all. The answer is that ``we all want to be seen as rational,'' and people confuse doubt with modesty and status, as they do with humility. I give no evidence for this. Displaying doubt does about as much good as wearing a lab coat.

Last, six rules. A rational doubt has a specific reason and exists to destroy its target; an unresolved or uninvestigated doubt does nothing; you may be proud of finishing off a cherished belief, but not of mere doubting; and to an ideal mind doubt would not be scary in the first place.
